%% Definisco la classe del documento: va specificato il formato della carta,
%% le dimensioni del carattere, di cominciare i nuovi capitoli a destra e
%% fronte-retro
\documentclass[a4paper,12pt,openright,twoside,titlepage]{book}

%% Imposto l'italiano
\usepackage[italian,english]{babel}

%% Imposto l'encoding dei caratteri. Specificando utf8 posso mettere
%% direttamente le accentate senza errori
\usepackage[utf8]{inputenc}
\usepackage[T1]{fontenc}

\usepackage{fancyhdr}

%% Formato del titolo dei capitoli
\usepackage[Lenny]{fncychap}
\ChNumVar{\fontsize{76}{80}\usefont{OT1}{pzc}{m}{n}\selectfont}
\ChTitleVar{\centering\Large\sffamily\bfseries}

%% Font (sceglierne uno)
\usepackage{mathpazo} % Use the Palatino font by default

%% Imposto l'indentazione per ogni paragrafo
\usepackage{indentfirst}

%% Lettera iniziale del capitolo su più righe
\usepackage{lettrine}

%% Interlinea a 1,5
\usepackage{setspace}
\onehalfspacing % interlinea 1.5
%\singlespacing %interlinea 1
%\doublespacing % interlinea 2

%% Questo per aforismi e citazioni
\usepackage{epigraph}

%% Per inserire gli url all'interno del testo
\usepackage[hidelinks]{hyperref}

%% Per inserire codici sorgenti
\usepackage{listings}
\usepackage{xcolor}
\definecolor{light-gray}{gray}{0.80}
\addto\captionsitalian{\renewcommand{\lstlistingname}{Esempio}}

%% Elenchi puntati e affini
\usepackage{enumitem}

%% Per includere le immagini: package e path dove andare a richiamarle per nome
\usepackage{graphicx}
\graphicspath{{./figures/}}

%% Questo serve per includere le immagini dove voglio io e non dove
%% pare all'editor!
\usepackage{float}

%% Fa sì che le tabelle possano essere divise tra più pagine
\usepackage{longtable}

%% Bibliografia
\usepackage[sorting=none,style=numeric-comp,backend=bibtex]{biblatex}
\addbibresource{biblio/bibliography.bib} % Database con le citazioni

%% Formattazioni varie
\pagestyle{fancy}\addtolength{\headwidth}{20pt}
\renewcommand{\chaptermark}[1]{\markboth{\thechapter.\ #1}{}}
\renewcommand{\sectionmark}[1]{\markright{\thesection \ #1}{}}
\cfoot{}

%%Testatine
%\rhead[\fancyplain{}{\bfseries\leftmark}]{\fancyplain{}{\bfseries\thepage}}
%\lhead[\fancyplain{}{\bfseries\thepage}]{\fancyplain{}{\bfseries\rightmark}}

%Pagine pari
\fancyhead[RE]{\bfseries\thesistitle}
\fancyhead[LE]{\bfseries\thepage}
\fancyhead[CE]{}
%\fancyfoot[LE]{\bfseries\thepage}
%\fancyfoot[CE]{}
%\fancyfoot[LE]{\textit{thesistitle}}

%Pagine dispari
\fancyhead[LO]{\bfseries\leftmark}
\fancyhead[RO]{\bfseries\thepage}
\fancyhead[CO]{}
%\fancyfoot[LO]{\bfseries\leftmark}
%\fancyfoot[CO]{}
%\fancyfoot[RO]{\textit{thesistitle}}
%\renewcommand{\footrulewidth}{0.4pt}

\fancypagestyle{plain}{%
	\fancyhf{}
	\fancyfoot[RO]{}
	\fancyfoot[LO]{}
	\fancyfoot[CO]{\bfseries\thepage}
	\renewcommand{\headrulewidth}{0pt}
	\renewcommand{\footrulewidth}{0pt}}

%%
%% Titolo della tesi, relatori et similia
%%
\newcommand{\thesistitle}{Blackbox Approach to CPS Reverse Engineering}
\newcommand{\supervisor}{Prof. \textbf{Massimo \textsc{Merro}}}
\newcommand{\cosupervisor}{Prof. \textbf{Severus \textsc{Piton}}}
\newcommand{\candidate}{\textbf{Marco \textsc{Oliani}}\\ Matr. VR457249}
\newcommand{\academicyear}{\large Academic Year 2022/2023}

%% Abstract
\newcommand{\fncyblank}{\fancyhf{}}
\newenvironment{abstract}%
{\cleardoublepage\fncyblank\null\vfill\begin{center}%
	\bfseries\abstractname\end{center}}%
{\vfill\null}

%######
%% Inizio documento
%######
\begin{document}

%###### Frontespizio ######%
\begin{titlepage}
	\begin{center}
		\thispagestyle{empty}
		
		% logo
		%\includegraphics[width=0.15\textwidth]{logo_univr.png}~\\[0.5cm]
		
		% Upper part of the page. The '~' is needed because \\
		% only works if a paragraph has started.
		
		\textsc{\LARGE Università degli Studi di Verona}\\
		%\textsc{Facolt\`a di Scienze Matematiche, Fisiche e Naturali}\\[1.0cm]
		\textsc{Department of COMPUTER SCIENCE}\\[1.0cm]
		\textsc{\large Master's Degree in \\ 
			COMPUTER SCIENCE AND ENGINEERING}\\[1.5cm]
		
		%\large Tesi di Laurea Magistrale \\[1.0cm]
		\large Master Thesis \\[1.0cm]
		% Title
		
		\hrulefill
		
		%{ \Large\bfseries [PROVVISORIO] \\
		%	BLACKBOX APPROACH TO \\ 
		%	CPS REVERSE ENGINEERING \\ }

		{ \Large\bfseries\MakeUppercase{\thesistitle}}
		
		\hrulefill \\[2.0cm]
		
		% Author and supervisor
		\begin{minipage}[t]{0.4\textwidth}
			\begin{flushleft} \large
				\emph{Supervisor:}\\
				\supervisor 
			\end{flushleft}
			\begin{flushleft} \large
				\emph{Co-supervisor:}\\
				\cosupervisor
			\end{flushleft}
		\end{minipage}
		\begin{minipage}[t]{0.4\textwidth}\raggedleft
			\begin{flushright} \large
				\emph{Candidate:} \\
				\candidate
			\end{flushright}
		\end{minipage}
		
		\vfill
		
		% Bottom of the page -> è la data
		\academicyear
		
	\end{center}
\end{titlepage}

\thispagestyle{empty}

%###### Aforisma ######%
\begin{titlepage}
	\thispagestyle{empty}
	\begin{flushright}
		\vspace*{40 mm}
		\small{\emph{``Gallina vecchia fa buon brodo''}\\(Valentino Rossi)}
	\end{flushright}	
	\vfill
\end{titlepage}
\thispagestyle{empty}

%###### ABSTRACT ######%
\begin{titlepage}
	\begin{abstract}
		Bla bla bla
	\end{abstract}
	\vfill
\end{titlepage}
\thispagestyle{empty}

%% Faccio partire la numerazione normale DOPO l'indice: lo stesso indice 
%% verrà numerato in numeri romani
\frontmatter
\tableofcontents
\mainmatter

%%%%%%%%%%%%%%%%%%%%%%%%%%%%%%%%%
%       Inizio TESI             %
%%%%%%%%%%%%%%%%%%%%%%%%%%%%%%%%%

%###### CAPITOLO 1 ######%
\chapter{Introduction}
\label{intro}

\lettrine[lines=2]{L}{orem} ipsum dolor bla bla bla. Ma dove metto l'abstract? Prova di interlinea che direi posso anche andare bene, ma bisogna poi vedere il tutto come si incastra alla fine, in modo da ottenere un bel risultato alla vista. Sperando di non esagerare con gli abbellimenti, per evitare che qualcuno leggendo si incazzi e mi faccia riscrivere tutto, che non ne ho minimamente voglia...

Lorem Ipsum is simply dummy text of the printing and typesetting industry. Lorem Ipsum has been the industry's standard dummy text ever since the 1500s, when an unknown printer took a galley of type and scrambled it to make a type specimen book. It has survived not only five centuries, but also the leap into electronic typesetting, remaining essentially unchanged. It was popularised in the 1960s with the release of Letraset sheets containing Lorem Ipsum passages, and more recently with desktop publishing software like Aldus PageMaker including versions of Lorem Ipsum.

Why do we use it?
It is a long established fact that a reader will be distracted by the readable content of a page when looking at its layout. The point of using Lorem Ipsum is that it has a more-or-less normal distribution of letters, as opposed to using 'Content here, content here', making it look like readable English. Many desktop publishing packages and web page editors now use Lorem Ipsum as their default model text, and a search for 'lorem ipsum' will uncover many web sites still in their infancy. Various versions have evolved over the years, sometimes by accident, sometimes on purpose (injected humour and the like).


Where does it come from?
Contrary to popular belief, Lorem Ipsum is not simply random text. It has roots in a piece of classical Latin literature from 45 BC, making it over 2000 years old. Richard McClintock, a Latin professor at Hampden-Sydney College in Virginia, looked up one of the more obscure Latin words, consectetur, from a Lorem Ipsum passage, and going through the cites of the word in classical literature, discovered the undoubtable source. Lorem Ipsum comes from sections 1.10.32 and 1.10.33 of "de Finibus Bonorum et Malorum" (The Extremes of Good and Evil) by Cicero, written in 45 BC. This book is a treatise on the theory of ethics, very popular during the Renaissance. The first line of Lorem Ipsum, "Lorem ipsum dolor sit amet..", comes from a line in section 1.10.32.

The standard chunk of Lorem Ipsum used since the 1500s is reproduced below for those interested. Sections 1.10.32 and 1.10.33 from "de Finibus Bonorum et Malorum" by Cicero are also reproduced in their exact original form, accompanied by English versions from the 1914 translation by H. Rackham.

Where can I get some?
There are many variations of passages of Lorem Ipsum available, but the majority have suffered alteration in some form, by injected humour, or randomised words which don't look even slightly believable. If you are going to use a passage of Lorem Ipsum, you need to be sure there isn't anything embarrassing hidden in the middle of text. All the Lorem Ipsum generators on the Internet tend to repeat predefined chunks as necessary, making this the first true generator on the Internet. It uses a dictionary of over 200 Latin words, combined with a handful of model sentence structures, to generate Lorem Ipsum which looks reasonable. The generated Lorem Ipsum is therefore always free from repetition, injected humour, or non-characteristic words etc.

Where can I get some?
There are many variations of passages of Lorem Ipsum available, but the majority have suffered alteration in some form, by injected humour, or randomised words which don't look even slightly believable. If you are going to use a passage of Lorem Ipsum, you need to be sure there isn't anything embarrassing hidden in the middle of text. All the Lorem Ipsum generators on the Internet tend to repeat predefined chunks as necessary, making this the first true generator on the Internet. It uses a dictionary of over 200 Latin words, combined with a handful of model sentence structures, to generate Lorem Ipsum which looks reasonable. The generated Lorem Ipsum is therefore always free from repetition, injected humour, or non-characteristic words etc.

%###### CAPITOLO 2 ######%
\chapter{Altro Capitolo}
Ri-bla \cite{einstein}

\chapter{Test include}
\label{test}

Lorem ipsum dolor e non mi ricordo più come va avanti. Che tanto manco è latino...


%###### BIBLIOGRAFIA ######%
\printbibliography[heading=bibintoc, title=Bibliography]

%% Vecchia versione
%\begin{thebibliography}{99}
%	\addcontentsline{toc}{chapter}{Bibliografia}	
%\end{thebibliography}

%%%%%%%%%%%%%%%%%%%%%%%%%%%%%%%%%
%		Fine TESI				%
%%%%%%%%%%%%%%%%%%%%%%%%%%%%%%%%%
\end{document}